% Generated draft: requires mathematical and source review.
\documentclass[11pt]{article}
\usepackage[margin=25mm]{geometry}
\usepackage{fontspec}
\setmainfont{Latin Modern Roman}
\usepackage{amsmath,amssymb,mathtools,unicode-math}
\setmathfont{Latin Modern Math}
\usepackage{enumitem,needspace,xurl,hyperref}
\setlength{\parindent}{0pt}
\setlength{\parskip}{5pt}
\setlength{\emergencystretch}{3em}
\newcommand{\N}{\mathbb{N}}
\newcommand{\Z}{\mathbb{Z}}
\newcommand{\Q}{\mathbb{Q}}
\newcommand{\R}{\mathbb{R}}
\newcommand{\C}{\mathbb{C}}
\newcommand{\F}{\mathbb{F}}
\newcommand{\A}{\mathbb{A}}
\newcommand{\PP}{\mathbb{P}}
\newcommand{\E}{\mathbb{E}}
\newcommand{\eps}{\varepsilon}
\newcommand{\dd}{\,\mathrm{d}}
\begin{document}
\section*{UnsolvedMath ID: 30005813}
\textbf{Source problem code:} catalogue entry
\subsection{Definitions and mathematical statement}
For a connected Riemannian manifold $(M,g)$, its diameter is
\[
\operatorname{diam}_g(M)=\sup\{d_g(x,y):x,y\in M\}.
\]
One possible meaning of the systole is the homotopy systole
\[
\operatorname{sys}_{\pi_1}(M,g)=\inf\{\operatorname{Length}_g(\gamma):\gamma\text{ is a piecewise smooth closed curve in }M\text{ that is not null-homotopic}\}.
\]
The supplied record does not specify whether this is the intended systole, rather than a homological systole or another convention. It also does not specify how manifolds with no noncontractible closed curves are treated. No convention for that case is imposed here.

Let $\mathcal C$ denote the intended class of Riemannian manifolds, which the supplied record does not fully specify. Once $\mathcal C$, the systole convention, and the meaning of topological type have been fixed so that the ratio below is defined, the supplied question has the following form:
\[
\text{For every }\varepsilon>0,\qquad
\#\left\{[M]_{\mathrm{homeo}}:\exists g\text{ with }(M,g)\in\mathcal C,\quad
\frac{\operatorname{sys}_g(M)}{\operatorname{diam}_g(M)}\geq 2-\varepsilon\right\}<\infty\;?
\]
Here $[M]_{\mathrm{homeo}}$ denotes the homeomorphism class of the underlying manifold; this makes explicit the usual interpretation of ``topological type,'' which still requires confirmation from the source. The metric is allowed to vary, and any proposed finite bound may depend on $\varepsilon$ and on the fixed parameters defining $\mathcal C$.

\textbf{Unresolved qualifications.} The record does not state whether dimension is fixed, whether manifolds are connected, compact, closed, complete, or permitted to have boundary, or whether curvature, fundamental-group, or other restrictions apply. For an ordinary real-valued ratio, the diameter must be positive and finite and the systole must have a specified finite value; treatment of other cases is missing. These are unresolved domain issues, not additional hypotheses supplied by this draft. The quantifier is exactly $\varepsilon>0$, with no stated upper bound. In particular, for $\varepsilon\geq2$ the displayed inequality imposes no restriction whenever the ratio is defined and nonnegative.

The corrupted original extract mentions a systolic bound involving $2$ and diameter and attributes to Gromov an equality characterization involving real projective space $\mathbb{RP}^{n}$. Its notation and hypotheses are incomplete, so neither the precise bound nor that characterization is asserted here as a verified theorem. The extract then runs into unrelated questions concerning Tverberg points, drawings of hypercube graphs, and orthogonal vectors; those passages do not specify the intended manifold class. Thus the displayed question is a conditional formalization of the supplied cleaned statement, pending resolution of its missing hypotheses.
\subsection{Short English statement}
For each $\varepsilon>0$, the supplied question asks whether only finitely many homeomorphism types admit a Riemannian metric whose systole-to-diameter ratio is at least $2-\varepsilon$. The intended manifold class, dimension restrictions, systole convention, and treatment of boundary or infinite quantities are unspecified, so a fully determined mathematical statement requires source review.
\subsection{Sources}
\begin{itemize}
\item Local UnsolvedMath registry, permanent ID \textbf{30005813}, title ``Topology Near the Systolic-Diameter Bound,'' source problem code \texttt{catalogue entry}. The supplied record belongs to the Oberwolfach Reports Open Problems collection and gives proposed year 2024; no proposer is identified.
\item Supplied report citation: \emph{Discrete Geometry} (2024), Oberwolfach Reports, DOI 10.4171/owr/2024/3, pp.~137--202. \url{https://doi.org/10.4171/owr/2024/3}. This is the authoritative report reference supplied by the record; no precise page locating this question is given.
\item Supplied literature source: \emph{Finiteness Theorems for Gromov-Hyperbolic Spaces and Groups}. \url{https://arxiv.org/abs/2109.13025}.
\item Literature assessment dated \textbf{2026-08-22}: the supplied curation labels the problem open, says that the report explicitly poses the systole-to-diameter finiteness question, and states that known finiteness theorems under stronger curvature, entropy, or hyperbolicity assumptions do not settle the stated $2-\varepsilon$ question. This is supplied dated curation, not an independently verified literature assessment.
\item The registry labels the statement \texttt{corrected\_verified} and reports that it was reconstructed and checked against the official Oberwolfach report. Those are supplied verification claims, not verification performed for this draft. The original extract attributes an incompletely stated equality characterization to Gromov without supplying a separate citation or its hypotheses.
\end{itemize}
\end{document}
